% Tracce uscite agli Esami di Stato (Sezione A, Settore Informazione) sugli argomenti del capitolo.
% Incluso da main.tex subito dopo il riassunto.

\section*{Tracce Uscite} % Usiamo * per non numerare la sezione nell'indice
\addcontentsline{toc}{section}{Tracce Uscite: Architettura del Calcolatore}

\subsection*{Prima sessione 2024 -- Elettronica \#2}
\textbf{Traccia}: Figura \#1 con il datapath di una semplice CPU RISC e le memorie esterne. (1) Unità funzionali (ALU, registri, PC\ldots); (2) stadi di esecuzione di un'istruzione e unità coinvolte in ciascuno; (3) funzionamento e vantaggi della pipeline; (4) problematiche dovute alla pipeline e come gestirle.

\textbf{Dove trovare la risposta}: sezione \emph{Tema Svolto: Datapath RISC e Pipeline (2024)}, che risponde ai quattro punti con riferimento alla figura; approfondimenti nelle sezioni \emph{Il Datapath}, \emph{La Pipeline} e \emph{Hazard}.

\subsection*{Prima sessione 2018 (seconda prova) -- Elettronica}
\textbf{Traccia}: struttura interna di un microcontrollore e alcuni esempi di sistemi elettronici in cui viene utilizzato.

\textbf{Dove trovare la risposta}: sezione \emph{Microcontrollori}, con la scaletta di risposta.

\subsection*{Prima sessione 2025 (Bologna) -- Ing. elettronica E03}
\textbf{Traccia}: consumo di potenza nei circuiti digitali CMOS; potenza statica e dinamica, cause della dissipazione, strategie di riduzione.

\textbf{Dove trovare la risposta}: sezione \emph{Potenza e consumo} (formula $P = \alpha C V^2 f$, leakage, clock gating, DVFS); per il transistore e l'inverter CMOS vedi il capitolo \emph{Elettronica}.

\subsection*{Prima sessione 2025 (Bologna) -- Ing. Informatica E05}
\textbf{Traccia}: chiamata a funzione e ruolo dello stack.

\textbf{Dove trovare la risposta}: capitolo \emph{Programmazione ad Oggetti e Algoritmi} e, per l'implementazione a livello di registri (\texttt{sp}, \texttt{ra}, argomenti), la sezione \emph{RISC-V: registri e formati} di questo capitolo.
